\documentclass[a4paper]{letter}
\usepackage[margin=22mm]{geometry}
\usepackage{csquotes}
\usepackage{siunitx}  % use for all units

% NOTE : Use as \si{\kph} or \si{\mps}
\def\kph{\kilo\meter\per\hour}
\def\mps{\meter\per\second}

\date{January 20, 2026}
\signature{The Authors}

\begin{document}
\begin{letter}{Ergonomics Review Response}
\opening{Dear Reviewer,}

Thank you for the careful review of our work. We revised the paper to address
your new comments.

\begin{displayquote}
  \textbf{Editor}: I have now received a further opinion from one of the
  original reviewers of your manuscript and comments are appended at the end of
  this email. You will see that although the reviewer judges there to be merit
  in the paper, s/he again advises that substantial revisions are required. The
  reviewer's comments appear reasonable and you are invited to consider these
  and modify your paper accordingly.
\end{displayquote}

Our responses and revision explanations are listed below the indented comments
of the reviewer. 

\begin{displayquote}
  \textbf{Reviewer 2}: The authors have made significant revisions. However,
  there are still issues need to be addressed, especially some statements are
  unrealistic, unreasonable, and unnecessary.
  \\\\
  Why do you provide a catalog? Have you seen a catalog on a journal article?
  Please delete P2.
\end{displayquote}

We assumed the journal reformats the article and undesired content from pages 1
and 2 would not be included. We have removed the table of contents from this
review draft.
 
\begin{displayquote}
  Tables. Please delete the unnecessary lines within a table.
\end{displayquote}

We have reduced the number of horizontal lines in the tables.

\begin{displayquote}
    P25. Ln 28. ``A stroller can be walked continuously for four hours only on
    tarmac, and only if the infant is lying in the cot.'' This is unrealistic.
    How could people walk and push an infant stroller for four hours??
\end{displayquote}

This sentence is referencing the four hour ISO~2631-1 threshold in relation to
our data (Figure~12). We have changed the sentence and the following one to:

``For tarmac, the vibration is relatively lower and longer durations are
possible. But for long durations, the parent/caregiver would need to consider
appropriate breaks to allow the infant to move.''

This reduces the chance it would be read as acceptable to stroll over tarmac for
four hours.

\begin{displayquote}
    5.2 Recommendations. Some of the recommendations, considering only the needs
    of infant strollers and ignore the needs of other road users, seem to be
    either unnecessary or impractical.
    For example, ``5. Designers and manufacturers should incorporate better
    suspension systems, …'' I think the designers and manufacturers have been
    trying to provide products with better suspension systems even without
    reading this paper because that is something intuitively reasonable.
    However, it probably depends on cost and price. Infant strollers are used
    for only a few months. So, people may be reluctant to pay an expensive price
    to get one with a sophisticate design considering vibration. It is common
    that parents push a stroller slowly and carefully so as to reduce the
    vibration.
    Another example is ``6. New cycling roads and sidewalks should have a
    tarmac-like smoothness…'' Roads and walkways are designed and constructed
    for multi-purposes and to fulfill the needs of the general population, not
    some specific users. Providing a smooth road surface could lead to other
    problems such as higher risk of slip/fall when the surfaces are wet. This is
    apparently undesirable for most road users.
    One more example, ``8. …..When transporting over tarmac, continuous
    durations should not exceed 4 hours,…..'' How could people transport an
    infant using a stroller or a cargo bicycle for four hours? This should
    definitely not occur. Even a healthy adult walking on a cycling road or
    sidewalks for four hours is very unlikely.
\end{displayquote}

We agree that a designer must consider a myriad of factors in their process to
arrive at an appropriate product or civil infrastructure. We have not researched
all factors that must be considered in the design of strollers, cargo bicycles,
infant seats, or roads in the scope of this paper. We stated ``based on the
findings of this paper'' in the opening statement of the recommendations section
for this reason.

To further reduce any confusion, we have changed the opening statement of the
recommendations to ``Nevertheless, the following list provides recommendations
for users, researchers, designers, and manufacturers based \textbf{only} on the
findings of this paper\textbf{, independent of other factors one must take into
consideration.} We believe these recommendations stand in spite of the
standard's limitations.'' to make this clearer. We also updated the
recommendation number 6, stating ``\textbf{To reduce road-induced vibration},
new cycling roads and sidewalks should have a tarmac-like smoothness.'' and
removed ``When transporting over tarmac, continuous durations should not exceed
4 hours, and need to include regular breaks, if only to allow the infant to move
free from the seat harness.'' from recommendation 8.

Regarding recommendation number 5, we provide evidence that typical modern
strollers lack adequate suspension. We have also had extensive communication
with stroller manufacturers in recent years. Some companies have even reviewed
this paper before we submitted it to Ergonomics. We have been surprised at the
lack of attention to vibration that seems to have been more present in past
designs. This recommendation is our reminder to users and companies not to
forget vibration, even when faced with a higher cost.

Regarding ``Providing a smooth road surface could lead to other problems such as
higher risk of slip/fall when the surfaces are wet.'', actually tarmac provides
adequate friction but may be regarded as not suitable or attractive for some
settings and may encourage speeding. We would rather stay away from such
discussions. 

Regarding ``transport of four hours'', we have personal experience with parents
transporting babies and children between cities on cycling tours that can be 4
to 8 hours in duration in a single day. Regarding ``This should definitely not
occur.'', we tend to agree but have not found any scientific evidence that
allows us to state a lower threshold. Our prior revision was formulated to
address this point:

``We did measure vibrations that would not be permitted for adults to maintain
long-term occupational health and it is reasonable to believe we should not
subject our infants to the same. The preceding statement should certainly not be
interpreted to mean that vibration lower than these limits is deemed acceptable
for infants. The relative vulnerability of infants points more towards infant
limits being lower than those of adults. Further research to establish direct
evidence to infant health and comfort would possibly permit more or less caution
than we conclude below in our recommendations.''

Given that our recommendations are based only on the findings in this paper and
this may have not been clear, we have left the recommendations the same other
than the addition to number 6 and deletion from number 8.

\closing{Sincerely,}
\end{letter}
\end{document}